%!TEX root = ../dissertation.tex

\chapter{Related Works}
\label{chp:relatedWorks}

In this section, I provide an overview of the related work on predictive process monitoring, from its role as a process mining methodology to its early approaches using flattened Object-Centric Event Logs (OCELs), followed by contemporary applications using heterogeneous graph structures, where I identify a key research gap, the absence of explainability analysis on heterogeneous predictive process monitoring for OCELs. Recognizing this gap, I examine methodologies on explainable artificial intelligence (XAI) used in other domains to provide an enhanced understanding of complex predictive models. Building on these insights I propose a four-component explainability framework to enhance existing predictive process monitoring pipelines by providing stakeholders with a transparent overview of the model's functionality.

\section{Process Mining}
    As data availability has increased exponentially so too has the need to use this data to provide actionable insights. Process Mining is a research discipline that seeks to act as a bridge between the fields of data mining and business process modeling and analysis. Its goal is to discover, monitor, and improve real processes by analyzing the events that make up a process. These events are presented in the form of an event log, a sequential record of events, such that every event refers to an activity (a well defined step in the process) and is related to a particular case or process execution \cite{vanderaalst2012}. 
    
    Aside from the sequence of events, event logs may also store additional information about each step in the process such as the person who performs the process, the timestamp of the event or a set of features related to elements tied to the event (such as the size of an order).
    
    Process mining can be broken down into three main types: process discovery, conformance checking, and process enhancement.
    Process discovery is the most prominent of the process mining techniques and refers to a methodology that transforms an event log into a model without any apriori information. This technique is crucial as it can show business stakeholders that the information present in an event log related to a process execution that may be further analyzed.
    Conformance checking, on the other hand, is a technique that validates whether a given model properly reflects the reality shown in the provided event log and vice versa. It is a diagnostics process meant to ensure the proper functioning of models created during process discovery.

    Finally, process enhancement techniques seek to extend or improve an existing process model using the information present in the event log. Their main goal is to take an encoding of features of the event data as input and deliver insights, predictions, or actions as output \cite{adams2022}. These insights can help identify problems in the underlying process, such as bottlenecks, or quantify important elements such as throughput times. Predictive process monitoring (PPM) methods stem from process enhancement techniques, their main goal being to provide a process's stakeholder with a prediction on the output of a currently running process execution. Through this prediction the stakeholder may be able to identify problems in the current execution and take preventive action in order to ensure the proper resolution of the process execution.

    \subsection{Object-centric processes vs. traditional processes}
        Traditional process mining techniques operate under two assumptions: Each object is of the same type and each event is associated with exactly one object. While these assumptions may be valid in certain cases, such as the handling of insurance claims (where each object describes an instantiation of an insurance claim) \cite{adams2022}, they do not properly represent most real-world processes. In real-life organizations, processes do not often happen in isolation, instead, several instances of different processes that are related by a set of objects can be executed simultaneously and interact with one another. In order to address this discrepancy a new paradigm was developed: object-centric processes. In object-centric processes an event may be related to objects of different types and an event with multiple objects may have multiple predecessor events \cite{gherissi2023}. These object-centric processes are represented through Object-Centric Event Logs (OCELs). OCELs store information on multiple aspects of a process through the the use of object types. For example, an OCEL representing an Order Management process will contain instances of different events in the process, such as the placing of the order or the creation of a package containing the ordered products, as well as information on the distinct objects that are related to each distinct event, such as the items ordered or the vehicle assigned for the package delivery. The resulting log's structure more closely resembles a graph rather than the sequential structure that was assumed for use in traditional process mining.

\section{Object-centric predictive process monitoring}
    Predictive process monitoring (PPM) on object-centric event logs (OCELs) refers to a process that monitors the running instances of a given object-centric process and provides predictions on key performance indicators (KPIs). The goal of these predictions is to alert stakeholders of possible risks arising in the process so they may mitigate their effects. 
        
    \subsection{From tabular to object-centric encodings}
        Process mining techniques, such as process enhancement, seek to provide insights by analyzing the information provided by event logs.
        Early process enhancement techniques encode features of event data in a tabular form where each row contains the feature values for an event. However, since each process execution is a temporally ordered sequence of events, summarizing the event data in a tabular form removes this temporal structure.
        Sequential encodings were then developed as a way to address this problem. By representing the data as a sequence of feature values the temporal structure of the information could be preserved. This, in turn, allowed for the use of specialized models that could leverage this sequential structure, such as LSTMs, to provide more accurate KPI predictions.

        While these methods are effective in analyzing and providing predictions on traditional event logs, object-centric processes and their OCELs provide a significant challenge due to their added complexity. 
        In order to apply traditional process enhancement techniques on OCELs initial implementations flattened the OCELs and encoded the object interactions as tabular features. 
        This initial approach was successful in showing the benefit the additional complexity found in OCELs could provide to predictive models. When combined with state of the art machine learning models such as Gradient Boosting predictions performed on object-centric processes showed better results than those trained using a naive approach \cite{galanti2023a}.
        However, while this method provides results the flattening process has major consequences on the quality of the calculated features. Flattening the OCEL can lead to events disappearing, a problem called deficiency. Conversely, events can also appear more than once, causing a problem called convergence. Finally, flattening can also cause the creation of misleading directly-follows relations, a problem labeled divergence. For example, an OCEL graph may show two events happening simultaneously but, when placed in a tabular form, the events are forced into an ordered state which creates a false directly-follows relation between one event and another.
        To combat this loss in feature quality Adams et al. proposed a graph-based encoding that retains the structural information of the process by maintaining the original graph structure of the OCEL \cite{adams2022}. The encoding also proposes an object-centric adaptation of features calculated for traditional event data \cite{deleoni2016}. The need for this translation arises from the two main differences between traditional event data and object-centric data: each event can have multiple predecessors/successors and each event might have multiple related objects of different types. All the feature calculations are adapted to the graph structure and are grouped according to different perspectives: Control-Flow, Data-Flow, Resource, Performance and Objects.


\section{Graph neural networks (GNNs)}
\label{sec:rw_gnns}
    The graph-based encoding proposed by Adams et al. can then be used to train Graph Neural Networks (GNNs) to provide a prediction on the chosen KPIs. 
    GNNs are special implementations of Neural Networks that are specifically created to handle graph based inputs and outputs. They're designed to leverage not just the features present in each node, but also the inherent relationship between objects in a graph. GNNs learn via a message-passing mechanism, built from two functions: an aggregation function and an update function. \cite{khemani2024}.
    In order to learn about the input structure the GNN model collects messages from nearby nodes to a particular node, which are then aggregated based on information gathered from each node's own neighborhood. This information is used to create node embeddings. The process continues recursively, with each iteration providing nodes with more information from their surrounding neighbors. Each iteration also provides the embeddings with information from distant parts of the graph. In the first iteration each node holds information from its immediate neighbors, nodes that are one hop away. In the next iteration, nodes aggregate information from nodes two hops away and updates each embedding. This process continues until each node's embedding includes data from its n-hop neighborhood. 
    This architecture enables GNNs to capture and learn from the complex relational structures present in graph data, making them effective for various tasks involving graph-structured data, such as graph classification, node classification, and link prediction \cite{khemani2024}. 
    Predictive process monitoring follows a classic supervised learning structure, where the GNN is provided with a tuple of features (in this case a graph representation of a current process) and an output value Y (the KPI value for the given process). Because of their ability to leverage the structural information present in the graph encoding of the OCEL GNNs have shown improvement in their predictive qualities when compared with a flattened encoding such as a tabular OCEL.
    Subsequent research has shown that by retaining the OCEL's structural information this graph-based encoding is able to outperform state of the art approaches using a flattened dataset \cite{adams2023}. 
    However, it's worth noting that a different work found the opposite to be true. A direct comparison between LSTM/Gradient Boosting models and GNN techniques across two different processes and 21 KPIs found that GNNs underperformed, while also taking considerably more time to train \cite{galantiDeleoni2024}.

    \subsection{A direct comparison baseline for heterogeneous GNNs}
        All the graph-based encodings mentioned so far worked with homogeneous graphs. These types of graphs are constructed under the assumption that all nodes represent the same object type. In the case of OCEL encoding, this approach requires aggregating object features in order to present the graph information as a series of interconnected event nodes. While graph encoding was able to properly address the problems caused by flattening the OCEL it still lacked the ability to properly encode all the distinct relationship types that existed in the event log between objects of different types and their associated events.
        In order to address this loss of information a heterogeneous graph structure was proposed that would treat each object instance as its own unique node type. Smit et al. \cite{smit2024} propose a framework for heterogeneous object-event graph encoding (HOEG) defined by two architectural choices: the use of a heterogeneous encoding that does not aggregate object features and the use of a k-dimensional heterogeneous GNN for its predictions.
        Under this encoding, it is no longer necessary to aggregate the distinct object features into each event node, instead, each object node holds its own set of distinct features directly related to that object. For example, a createOrder event no longer needs to aggregate the object information for the order and any other related object such as the items that make up the order. In a heterogeneous encoding each of those object types is given its own separate node with its own distinct features.
        Smit et al.'s work showed that expanding the representative power of the graph led to inconclusive results in the predictive quality of their models. While the heterogeneous structure provided better predictions on one of the tested datasets, it tied on a different one, and provided lower results on another.
        My thesis adopts their work as its direct comparison baseline in the Models chapter, however it is necessary to address some key differences. First, Smit et al.'s k-dimensional GNN is a homogeneous GraphConv-based operator, made heterogeneous through PyTorch Geometric's generic to\textunderscore hetero() transform. Second, their graph structures have no object-to-object edges, presenting only the event-to-event and object-to-event relations. Finally, HOEG predicts remaining time at every event node of a complete process execution in a single forward pass (no masking to one target), instead of once per prefix. For my work I chose to utilize a heterogeneous model architecture based around HGT convolutions which are natively type-aware and require no model wrapper, the graph structure used to train and test this model contains all three types of node connections (object-to-event, event-to-event and object-to-object), and the training/testing sets are composed of distinct prefix subgraphs drawn from every process execution in the dataset.

\section{Explainable AI}
\label{sec:rw_xai}
    Explainable AI (XAI) refers to the field of study of methodologies used for providing end users with an understanding of the way complex machine learning models function.
    
     \subsection{Explainability on Neural Networks}
     \label{sec:rw_exp_on_gnn}
        Given the widespread availability of structured data more and more business are adopting predictive models for high-stake decision-making processes. Simple predictive models such as decision trees are able to generate their own explanations and are referred to as white box methods. Their methodology is easily interpretable and an end user is able to understand why the model predicted a certain value by following the logic within the model. However, their simplicity also limits their overall predictive power and range of applications. 
        Since the concrete goal of these models is to provide as accurate predictions as possible the current trend is to increase the predictive performance with the use of deep learning methods such as neural networks. While these methods offer a powerful and versatile solution to predictive problems their inherent complexity leads to a lack of transparency and an inability to interpret the given predictions. Identified as black box methods, these interpretability problems can lead to them not being used for high-stake decision-making processes even with their superior performance. 
        
        The explainable artificial intelligence (XAI) paradigm seeks to address this issue by providing tools to help make these black box methods more understandable to the end user. XAI methods try to approximate the behavior of the model
        with post-hoc explainable techniques such as Shapley values, feature importance, attention weights, etc. 
        While the concepts of explainability and interpretability are often used interchangeably a subtle distinction must be highlighted \cite{stevens2022}. The first one refers to the ability to properly link the effect an input has on a model's output, i.e. how faithful the explanation is to the model. The second refers to the ability to express this relationship in an unambiguous and simple way so that an end user may properly use the information. Therefore, a simple explanation may be easily interpretable but may not be faithful to the predictive model. Therefore, it is especially important to strike a balance within the two, aiming to provide explanations that are faithful to the model but also easily interpretable.
        
        There exists a variety of XAI methodologies for providing interpretable explanations of a given model. Three widely used methods rely on perturbation (sometimes also referred to as ablation), feature attribution and counterfactual analysis.
        Perturbation methods seek to quantify the impact of removing or altering input values and measuring the effect this change has on the model's prediction. An example of these methods are Shapley Values. This explainability method adapts the game theory approach of Shapley Values, which seeks to find a fair distribution of a game's payout among the players that have collaborated. Under this method, the features of an instance correspond to the players and the payout is the difference between the prediction made by the predictive model and the average prediction. Thus, for each instance, the Shapley value of a feature provides a quantifiable expression of its impact on the model's prediction \cite{lundberg2017}. 
        Feature attribution methods seek to clarify the roles of input features in the prediction of the outcome and have been successfully used in different fields such as molecular science and Internet of Things projects.
        Counterfactual explanations seek to identify the minimal changes needed in an input in order to change the model's prediction. These methods have found success in domains such as recommendation systems and molecular science.

        \subsection{Explainability on predictive process monitoring}
        The proper application of explainability methods is especially important for predictive process monitoring since the adoption of black box models in critical business areas is tied to the trust stakeholders can place on the predictions being provided by the model. As Galanti et al. \cite{galanti2023b} show in their work on explainable support systems, process analysts and stakeholders will not adopt predictive models that do not put explanations as a core feature and build the user's trust. Because of this imperative, much work has been done to extend predictive process monitoring with explainability features. Mehdiyev et al. \cite{mehdiyev2025} provide a comprehensive survey of the existing work in this field. While their review shows a large amount of PPM applications using a variety of XAI methods the majority of the works reviewed adhere to the traditional process mining paradigm. The authors distinctly point out the need for further work in the field of object-centric process explainability, highlighting that the change in paradigm challenges traditional explanation approaches. 
        
        Further research into works focused on object-centric processes reveals the following. Weinzier et al.\cite{weinzierl2020}, Wickramanayake et al.\cite{wickramanayake2022}, and Pasquadibisceglie et al. \cite{pasquadibisceglie2021} seek to explain the object-centric processes by focusing directly on the Deep Learning methods chosen. Their works provide explanations through the use of Layer-wise relevance propagation, attention mechanisms, and neuro fuzzy networks respectively. These methods show positive results but are limited by being model specific approaches that are specifically designed for certain model types. Model-agnostic approaches include Galanti et al. \cite{galanti2023b} \cite{galanti2023a} who adopted the use of Shapley Values as an explainability method for their object-centric process analysis, as well as, Hsieh et al.\cite{hsieh2021} and Huang et al.\cite{huang2022} who propose the use of counterfactual methods to provide explanations. While all these approaches, both model-centric and model-agnostic, successfully delivered interpretable explanations they were applied to tabular encodings of object-centric data and relied on traditional neural network implementations. On the other hand, works such as Adams et al.\cite{adams2022}, Gherissi et al. \cite{gherissi2025}, Smit et al.\cite{smit2024} and de Leoni \& Volpato \cite{deleonieVolpato} focus on the use of graph-based encodings and GNNs for  generating predictions but do not offer any type of explainability method for interpreting their models' results. As far as I could tell, no work has been done to apply explainability techniques on object-centric processes encoded as heterogeneous graphs.
    
        \subsection{The GNN-explainability landscape}
        Since current state of the art PPM methods are based around the use of GNNs it is essential to explore XAI methods specifically designed for graph-based models. Broad taxonomic surveys \cite{nandan2025} have been conducted to map out the universe of options available in the space. 
        Two main explanation categories exist: Factual and Counterfactual. Factual explanations aim to explain a model by identifying the input features that are the most impactful on the model's prediction. Counterfactual explanations, on the other hand, seek to explain the model by identifying the minimum change to the input graph that would lead to a different prediction. Counterfactual explanations can be used to identify a set of related characteristics whose modifications can alter the model's prediction.
        Factual explanations can be further categorized as either post-hoc or self-interpretable. Self-interpretable techniques seek to incorporate interpretability directly into the architecture and training process of the GNN. Such explanations tend to employ graph attention techniques to assign attention weights to different nodes and edges, using these to identify the most prominent elements in the input graph structure.
        Post-hoc techniques, on the other hand, treat pre-trained GNNs as black boxes. These techniques can be broken down into instance-level approaches and model-level, according to their need to access the internal parameters of the model.
        Instance-level approaches include two further categorizations (among several) which were used in this thesis implementation: perturbation-based methods and gradients/feature-based methods.

    \subsubsection{Perturbation and gradient/feature-based methods}
    \label{sec:rw_pert_exp}
        Perturbation-based methods seek to determine the characteristics or traits that are critical for the model to provide accurate predictions by manipulating the model's inputs and observing the changes caused in the output. Perturbation-based methods provide a straightforward way of understanding the function of elements in models.

        Since GNNs are trained to obtain information from an input's nodes as well as its structure it is important that perturbations respect this structure as much as possible. In GNNs these perturbations are applied to the model in the form of masks, which limit the available information in the input that is passed to the predictive model without significantly altering the input's actual structure. Methods such as Shapley Values, mentioned above, or Leave-One-Out (LOO) methodologies \cite{limonroejurafsky2017}, which apply the same logic as Shapley values but consider only one feature at a time as opposed to the entire feature space, can be extended for use with GNNs by applying this mask methodology. The main setback of Shapley Values is their computational intensity, having to obtain the exact coalition enumeration for a complex graph structure can be extremely costly. However, this cost can be mitigated through the use of estimation methods such as Monte Carlo permutation sampling described in Castro et al.'s ApproShapley algorithm \cite{castro2009}\cite{lundberg2017}.

        Masks can also be directly optimized via backpropagation, allowing for methods that identify a mask which shows only the most relevant characteristics for an accurate prediction at a lower computational expense. The resulting masked version of the original graph constitutes a subgraph known as the explanation subgraph. 
        
        GNNExplainer is a perturbation method that utilizes masks to learn important patterns in both edge and node features. It is a model-agnostic implementation and identifies a succinct subgraph that is most relevant or influential for a GNN's prediction of a particular instance \cite{ying2019}. While the original GNNExplainer implementation was designed for homogeneous models it can be adapted for use with heterogeneous models by creating a separate mask per node type. However, a limiting factor of this implementation is that GNNExplainer is not able to perform edge-masking and cannot provide edge importance values for heterogeneous graphs. 
        
        PGExplainer is another perturbation based explanation method that utilizes masks to change the model's input graph. However, PGExplainer parametrizes mask generation across instances via a trained neural network instead of optimizing per-instance in order to improve generalization \cite{luo2020}. 
        
        While both model options were successfully tested on the thesis dataset PGExplainer was decided against because its implementation required a greater degree of processing to properly integrate into the heterogeneous framework. GNNExplainer is retained only as a standalone comparison point against the chosen implementation since its lack of edge importance severely limits its application and its assessment metrics were consistently lower than those of the chosen explanation methodology.  
        Instead, the final perturbation-based explanation works on a combination of exhaustive LOO and Shapley Values. Exhaustive LOO analysis works as a cost efficient way to identify the most impactful elements in the input's structure by running on every feature of the input graph (nodes, edges and node features) and identifying those with the greatest impact on the prediction value. By limiting the explanation space to the elements identified through LOO and using Monte Carlo sampling I was able to use Shapley Values to properly quantify the impact those graph features have on the model's prediction without incurring major computational cost.
        
        Gradients/feature-based methods identify key input features by calculating the gradients of the target prediction with respect to those features. These methods were not originally designed with GNNs in mind but, given their simplicity and generality, they were expanded into the graph domain \cite{sundararajan2017}. These methods can be used to provide a computationally efficient overview of the most important features present in the nodes that make up a graph. In my framework the gradient/feature based explanation is used to identify the most important features in the previously identified most important nodes and quantify their importance for the predictions through attribution values. 

    
    \subsubsection{Model-level explanations}
        Model-level explanations provide explanations for a GNN that are not dependent on any particular input sample. They aim to provide a more general understanding of the model and its ability to elucidate generic behaviors. GNNInterpreter is an example of a model-level explainer that was considered for use in the current project. It provides explanations by learning a probabilistic graph distribution that reveals a class's discriminative pattern rather than explaining single predictions \cite{wang2023}. 
        Both self-interpretable and model-level explanations were considered for this work but were ultimately left out. 
        Given that one of the goals of the thesis was to provide a general framework for applying explainability to PPM, post-hoc explanations can ensure the explainability block will work well with any given prediction model and doesn't need to rely on the explanation method being a part of the model itself. Model-level explanations were left out because there were no proper implementations found for working with heterogeneous graph neural networks using these techniques and an extension of these methods for this particular use was outside the scope of the work \cite{yuan2022} \cite{nandan2025}.

    \subsubsection{Counterfactual Explanations}
        Counterfactual techniques aim to explain a model by identifying the minimal change that could be performed on the input graph to generate a change in the model's prediction. They offer a good complement to factual explanations by expanding a user's understanding of why a particular prediction is being made with an explanation of why a prediction was made instead of a different one  \cite{pradoromero2024}.
        Like factual explanations, counterfactual explanations can also be separated into instance and model-level explanations, my work focuses entirely on instance level.
        Counterfactual explanations can further be broken down into three main groups: perturbation-based, heuristic-based, and search-based.
        Perturbation-based techniques, much like their factual counterparts, consist of modifying the existing graph by adding or eliminating edges in the graph until a change in the prediction made by the model is achieved. The resulting graph is considered a counterfactual of the original input.
        Heuristic-based methods rely on a rigorous policy of modifying the input graph until they reach a valid counterfactual.
        Search-based techniques rely on using search methods within the counterfactual space to identify relevant candidates. The main challenge of this technique is either properly limiting the search space or employing sufficiently efficient search algorithms that the exponential search space is not a problem. An example of this methodology called NICE offers a general nearest-unlike-neighbor counterfactual algorithm for heterogeneous tabular data \cite{brughmans2022}.
        Another example of this methodology used in PPM is the LORELEY method which extends LORE (a counterfactual method that generates explanations using genetic algorithms) with control-flow constraints for realistic counterfactuals in predictive business process monitoring \cite{huang2022}.

    \subsubsection{Explainability metrics}
    \label{sec:rw_exp_metrics}
        Given the width of available XAI options that exist it was also necessary to provide a way to provide a comprehensive evaluation of their effectiveness. 
        While ground-truth metrics can be used to calculate an explainer's accuracy on a synthetic dataset real-world applications will rarely provide such information. To this end, metrics that can be used to assess a model's technique in the absence of a ground-truth were developed \cite{stevens2022}. The assessment on an explainer's technique aims to evaluate how good the explainer is at generating explanations that humans can comprehend. Metrics such as fidelity, a measure of the impact the important nodes the explainer identified have on the model's output, and sparsity, which quantifies the proportion of features that are deemed significant by the explanation, can provide a concrete idea of the effectiveness of the chosen explainer \cite{yuan2022}.

    \subsubsection{Explainability for heterogeneous graphs}
        While there are many XAI options for use with GNNs the vast majority of the existing methods have no direct implementation for heterogeneous graph neural networks, focusing only on homogeneous graph structures \cite{yuan2022} \cite{nandan2025}. This severely limits the available options for the current work and emphasizes the need for developing specialized explainability techniques that can handle the complexity of Heterogeneous Graph Neural Networks.
        Aside from the chosen implementation, two other different methodologies were considered. The first was an implementation of HGExplainer. This explanation method extends the GNNExplainer mask/mutual-information paradigm to heterogeneous graphs via meta-path sampling. However, its lack of open-source implementation and its proposed scope being specifically for link-prediction/recommender-systems tasks, not general regression made it infeasible to apply directly to the project \cite{mika2023}.
        A second option that was investigated was HTGExplainer. This method extends the same mask/mutual-information paradigm seen in previous explanation methods to heterogeneous and temporal GNNs making it an ideal fit for PPMs. The model has a publicly available implementation which was investigated and confirmed to work. However, the original work is based on a DGL implementation, a port was attempted to this project's chosen framework of Pytorch Geometric but it was unsuccessful \cite{li2023}. 

        As discussed further below, the final chosen methodology is based on the work of Zhai et al.\cite{zhai2025}, however certain aspects of their implementation were changed to better suit the data and analysis of this thesis. While their work adapts explainability methods for heterogeneous structures their model relies on a heterogeneous graph with only 1 node type, where distinct nodes are identified by a node feature as opposed to being modeled as distinct graph objects. 
        Given this difference, the ablation analysis was replaced by a perturbation-based explanation combining exhaustive LOO to identify the key elements in the structure and Shapley Values to quantify their impact. Gradient/feature-based explanations are adopted as a second component to identify the important features within the most important nodes in the graph. The counterfactual implementation is adapted taking this structural difference in mind. Additionally a fourth component is added to the explainability layer in the form of metrics to assess the effectiveness of the explanations provided.

        Table \ref{tab:xai_comparison} summarizes the methods surveyed in this section across the dimensions most relevant to this thesis's setting. The table indicates whether a method natively supports heterogeneous node types, whether it can attribute importance to edges (not just nodes), whether it offers factual and/or counterfactual explanations, whether it has been applied to temporal/process data, whether an open-source implementation exists, and how each was ultimately treated in this project.

        \subsubsection{Methodological template}
        Of the eight methods surveyed, Zhai et al.'s \cite{zhai2025} is the only one that has been applied to object-centric predictive process monitoring specifically. Additionally, it offers open-source availability and coverage of both factual and counterfactual explanations. The remaining methods, on the other hand, target only general graph tasks, non-object-centric predictive process monitoring, or non-graph domains. Because of this, the explainability structure set forth by Zhai et al. \cite{zhai2025} in order to provide post-hoc explanations on a heterogeneous GNN framework was chosen as this thesis template. This methodology provides end users with explanations by focusing on three types of XAI methods: ablation analysis, feature attribution methods and counterfactual explanations. Their explanation methodology provides a holistic view of the information provided to the model by quantifying the importance the model is assigning to each feature and the changes that could be made to a given input in order to change the output. 


        \begin{landscape}
        \begin{table}[t!]
        \centering
        \small
        \renewcommand{\arraystretch}{1.15}
        \caption{Comparison of the GNN-explainability methods surveyed in this section.}
        \label{tab:xai_comparison}
        \begin{tabulary}{\linewidth}{>{\raggedright\arraybackslash}p{2.3cm}LLLLL>{\raggedright\arraybackslash}p{4.5cm}}
            \toprule
            Method & Node types & Edge importance & Explanation type & Temporal & Open-source & Disposition in this thesis \\
            \midrule
            GNNExplainer \cite{ying2019} & Heterogeneous via per-type mask workaround & $\times$ on heterogeneous graphs & Factual & $\times$ & \checkmark & Tested; retained as comparison-only baseline \\
            PGExplainer \cite{luo2020} & Homogeneous & Not evaluated & Factual & $\times$ & \checkmark & Tested; rejected -- heterogeneous integration too costly \\
            GNNInterpreter \cite{wang2023} & No heterogeneous implementation found & --- (model-level) & Factual, model-level & $\times$ & Not confirmed & Reviewed only; rejected -- model-level scope, no heterogeneous support \\
            NICE \cite{brughmans2022} & N/A -- tabular, not graph & --- & Counterfactual & --- & Not stated & Reviewed only; background context \\
            LORELEY \cite{huang2022} & N/A -- sequence-based & --- & Counterfactual (+ factual rules) & Implicit (sequential) & Not stated & Adapted as this thesis's counterfactual template, in a regression setting with graph dissimilarity instead of Levenshtein distance \\
            HGExplainer \cite{mika2023} & Heterogeneous (meta-path) & Not confirmed & Factual & $\times$ & $\times$ & Reviewed only; rejected -- no open-source implementation, link-prediction/recommender-systems scope only \\
            HTGExplainer \cite{li2023} & Heterogeneous & Not confirmed & Factual & \checkmark\ (only method explicitly temporal-capable) & \checkmark, DGL only & Tested in DGL, confirmed working; rejected -- PyTorch Geometric port unsuccessful \\
            \textbf{Zhai et al. \cite{zhai2025}} & Heterogeneous architecture; single node type in their own study & Partial -- factors into counterfactual dissimilarity, not per-instance masks & Both (factual + counterfactual + ablation) & $\times$ & \checkmark, vendored & \textbf{Chosen as this thesis's template}, extended to true multi-node-type graphs \\
            \bottomrule
        \end{tabulary}
        \end{table}
        \end{landscape}
        
\section{Summary and proposed framework}
    In summary, given the identified gaps in explainability research for predictive process monitoring I propose an explainability framework that adapts and extends the work set forth by Zhai et al. The framework is composed of four components: a perturbation-based explanation, a gradient/feature-based explanation, a counterfactual explanation and a set of metrics to assess the explanation provided to the end user. This framework is designed to be model-agnostic and expand on current state of the art predictive process monitoring methods that use heterogeneous Graph Neural Networks to provide KPI estimations.